\scanpage{34}
\deleted{\deleted{\unsure{Soit|Théorème}} Soit $F$ un faisceau de $l$-torsion sur $Y$, et soit $\xi\unsure{\in|\,} H^n(X,f^*(F))$. Pour que $\xi\in\mathrm{Im}\ H^n(Y,F)$, il f. et s. que pour \deleted{tout} $y\in Y$, la classe $\xi_{\bar y}\in H^n(X_{\bar y},F_{\bar y})$ \deleted{\unsure{induit|induite}} sur la fibre géométrique de $X_{\bar y}$ soit nulle.}

\margin{\deleted{\unsure{IX|X}}}

\deleted{\uline{Corollaire} Supposons $f$ \inserted{qu. cp et} localemt $(n,\unsure{\,|\ell})$-acyclique pour $l$. Pour qu'il soit \uline{globalemt} $n$-acyclique pour $l$, il f. et s. que ses fibres géométriques soient $n$-acycliques pour $l$.}

\margin{$Y$ loc. noeth.}

\deleted{Dém. du th. 2.}

\deleted{Par hypothèse on a}
\[
\deleted{\left\{\begin{aligned}
R^q f_*(f^*(F))&=0 &&\text{pour }0<q\le n-1\\
R^{\unsure{0|q}}f_*(f^*(G))&\simeq G
\end{aligned}\right.}
\]

\deleted{donc la \deleted{suite} spectrale donne une suite exacte [NB. on suppose $n-1\ge0$ i.e., $n\ge1$\,; si $n=0$, \unsure{et|est} \unsure{mérite|nécessite} un énoncé séparé avec le \unsure{gpe|prop} fond.]}
\[
\deleted{\unsure{-|\longrightarrow|\cdots\longrightarrow} H^{\unsure{n}}(Y,F)\to H^n(X,F)\to H^0(Y,R^n f_*f^*(F))}
\]

\deleted{donc on est ramené à prouver qu'une section de $R^n f_*(f^*(F))$ qui induit 0 sur chaque fibre géom. est nulle. On est ramené}

% archive sheet number: 33
